\documentclass[12pt,a4paper]{article}

\usepackage[utf8]{inputenc}
\usepackage[T1]{fontenc}
\usepackage[ngerman]{babel}
\usepackage{amsmath,amssymb,amsthm,mathtools}
\usepackage[colorlinks=true,linkcolor=blue,citecolor=red,urlcolor=blue]{hyperref}
\usepackage{geometry}
\geometry{margin=2.5cm}
\usepackage{enumitem}
\usepackage{booktabs}
\usepackage[capitalise,noabbrev]{cleveref}

\newtheorem{theorem}{Theorem}[section]
\newtheorem{lemma}[theorem]{Lemma}
\newtheorem{proposition}[theorem]{Proposition}
\newtheorem{corollary}[theorem]{Korollar}
\newtheorem{conjecture}[theorem]{Vermutung}
\theoremstyle{definition}
\newtheorem{definition}[theorem]{Definition}
\theoremstyle{remark}
\newtheorem{remark}[theorem]{Bemerkung}

\newcommand{\Ksym}{K_\sigma^{\mathrm{sym}}}
\newcommand{\EE}{\mathbb{E}}
\newcommand{\muG}{\mu_\sigma}
\newcommand{\PP}{\mathcal{P}}

\title{Gerade Dominanz der quadratischen Weil-Form:\\ Spektralanalyse, computergestützter Beweis\\ und das Resolventen-Gap-Theorem\\[6pt] \large\textit{Kurzfassung -- vollständige Beweise in der englischen Version}}
\author{Lukas Geiger\\
\small Bernau, Deutschland\\
\small ORCID: \href{https://orcid.org/0009-0005-7296-1534}{0009-0005-7296-1534}}
\date{\today}

\begin{document}
\maketitle

\begin{abstract}
Dies ist die zweite Arbeit in einer dreiteiligen Serie (Teil~I: Grundlagen und Hindernisse \cite{Geiger2026partI}; Teil~III: Conclusio \cite{Geiger2026partIII}). Wir befassen uns mit dem offenen Problem, das in Connes' Übersicht von 2026 \cite{Connes2025} identifiziert wurde: ob der kleinste Eigenwert der quadratischen Weil-Form $\mathrm{QW}_\lambda$ einfach ist und eine gerade Eigenfunktion besitzt. Wir etablieren drei Hauptergebnisse. Erstens das \emph{Shift-Parity-Lemma}: Die Differenzmatrix zwischen den Operatoren des geraden und ungeraden Sektors hat für alle $N \geq 3$ Moden und alle normierten Verschiebungen $r \in (0,2)$ einen streng negativen kleinsten Eigenwert, was impliziert, dass jede Primzahl individuell gerade Eigenfunktionen bevorzugt. Zweitens zertifizieren wir die gerade Dominanz rigoros durch einen computergestützten Beweis unter Verwendung von Intervallarithmetik für $33$ Werte von $\lambda$ zwischen $100$ und $1{,}300{,}000$, wobei die zertifizierten Lücken (Gaps) monoton wie $\sim\!\sqrt{\lambda}$ wachsen. Drittens beweisen wir das \emph{Leading-Mode-Cancellation-Lemma}: Die Überlappungsdifferenzen zwischen dem geraden und ungeraden Sektor heben sich paarweise mit der exakten Konstanten $c = 2 + \sqrt{2}$ auf, was $(E_{\sin} - E_{\cos})/D(\lambda) = O(1/\log\lambda) \to 0$ ergibt. Dieser Resolventen-gedämpfte Rahmen reduziert den verbleibenden analytischen Schritt (M1$''$) auf klassische Fourier-Analyse und den Primzahlsatz.
\end{abstract}

\noindent\textbf{MSC 2020:} 11M26, 11M36, 47A75, 65G20\\
\textbf{Schlagworte:} Riemannsche Vermutung, quadratische Weil-Form, gerade Dominanz, Shift-Parity-Lemma, computergestützter Beweis

\newpage
\tableofcontents

% =============================================================================
\section{Einleitung}
\label{sec:intro}
% =============================================================================

Diese Arbeit ist die zweite in einer dreiteiligen Serie, die die Riemannsche Vermutung durch Eichtheorie, Spektralanalyse und arithmetische Thermodynamik untersucht. Teil~I \cite{Geiger2026partI} etabliert die thermodynamische Landschaft (strukturelle Ergebnisse R1--R9) und dokumentiert das Scheitern des eich-theoretischen Ansatzes (Hindernisse K1--K4), was zu einer Neuausrichtung auf das Spektralprogramm von Connes führte. Teil~III \cite{Geiger2026partIII} liefert die Conclusio.

Hier entwickeln wir den Spektralansatz vollständig. Gemäß Connes' Übersicht von 2026 \cite{Connes2025} untersuchen wir die quadratische Weil-Form $\mathrm{QW}_\lambda$ und ihre Gerade-Ungerade-Zerlegung. Unsere drei Hauptbeiträge sind:
\begin{enumerate}
  \item Das \textbf{Shift-Parity-Lemma} (\Cref{sec:shift-parity}): Ein rein algebraisches Ergebnis, das zeigt, dass jede Primzahl einzeln gerade Eigenfunktionen bevorzugt, bewiesen über trigonometrische Einträge in geschlossener Form.
  \item \textbf{Computergestützte Zertifizierung} (\Cref{sec:certificates}): Rigorose Verifizierung der geraden Dominanz für $33$ Werte von $\lambda$ zwischen $100$ und $1{,}300{,}000$ mittels Intervallarithmetik, wobei die zertifizierten Gaps wie $\sim\!\sqrt{\lambda}$ wachsen.
  \item Das \textbf{Leading-Mode-Cancellation-Lemma} (\Cref{lem:leading-cancel}): Die Überlappungsdifferenzen zwischen den Sektoren heben sich paarweise mit der exakten Konstanten $c = 2 + \sqrt{2}$ auf, was das analytische Gap (M1$''$) auf die Fourier-Analyse und den Primzahlsatz reduziert.
\end{enumerate}


% =============================================================================
\section{Connes' Reduktion und die quadratische Weil-Form}
\label{sec:connes}
% =============================================================================

Eine aktuelle Übersicht von Connes \cite{Connes2025} reduziert die Riemannsche Vermutung auf eine spektrale Bedingung der quadratischen Weil-Form $\mathrm{QW}_\lambda$. Wir fassen den Zusammenhang zusammen und präsentieren numerische Belege für den verbleibenden offenen Schritt.

\subsection{Connes' Reduktion}

Connes konstruiert einen selbstadjungierten Operator $A_\lambda$ auf $L^2([-\log\lambda,\,\log\lambda])$ aus der expliziten Formel von Weil. In additiven Koordinaten ($u = \log x$) lautet die quadratische Form
\begin{multline*}
  \langle A_\lambda f \,|\, f \rangle = (\log 4\pi + \gamma)\|f\|^2 \\
  + \int_0^\infty \Big[\langle T_u f + T_{-u} f,\, f \rangle - 2e^{-u/2}\|f\|^2\Big]\,\frac{e^{u/2}}{2\sinh(u)}\,du \\
  + \sum_p \sum_{m=1}^\infty \frac{\log p}{p^{m/2}} \langle T_{m\log p} f + T_{-m\log p} f,\, f \rangle,
\end{multline*}
wobei $T_u$ den Shift-Operator bezeichnet, $\gamma$ die Euler-Mascheroni-Konstante ist und der archimedische Kern $e^{u/2}/(2\sinh u)$ aus der Variablen-Substitution $x = e^u$ resultiert, angewendet auf Connes' Gleichung~(10) in \cite{Connes2025}, mit $x^{1/2}/(x-x^{-1})\,d^*\!x = e^{u/2}/(2\sinh u)\,du$. Der Operator $A_\lambda$ ist nach unten beschränkt mit kompaktem Resolventen (daher diskretes Spektrum).

\begin{theorem}[Connes, Theorem 6.1 in \cite{Connes2025}]
Sei $L > 0$ und $D$ eine reelle Distribution auf $[0,L]$ mit gerader Fortsetzung $\tilde{D}$ auf $[-L,L]$. Angenommen, die quadratische Form mit Kern $\tilde{D}(x-y)$ definiert einen nach unten beschränkten selbstadjungierten Operator auf $L^2([-L/2,\,L/2])$, dessen kleinster Eigenwert einfach, isoliert und mit einer geraden Eigenfunktion $\eta$ ist. Dann liegen alle Nullstellen von $\tilde{\eta}(z)$ auf der reellen Geraden.
\end{theorem}

\noindent Connes' Abschnitt 6.6 stellt fest: „Um Theorem 6.1 anzuwenden, muss man zeigen, dass der kleinste Eigenwert von $\mathrm{QW}_\lambda$ \emph{einfach} mit einem \emph{geraden} Eigenvektor ist.“ Dies ist der verbleibende offene Schritt.

\subsection{Galerkin-Diskretisierung}

Wir approximieren $A_\lambda$ in der Cosinus-Basis $\{\varphi_n\}_{n=0}^{N-1}$ (gerader Sektor) und der Sinus-Basis (ungerader Sektor) auf $[-L,L]$ mit $L = \log\lambda$. Die Matrixelemente werden mittels Quadratur berechnet ($n_{\mathrm{quad}} \geq 800$, $n_{\mathrm{int}} \geq 500$) unter Verwendung des korrigierten Kerns $e^{u/2}/(2\sinh u)$ und Primzahlen $p \leq \max(\lambda, 100)$.

\subsection{Ergebnis 1: Spektrales Gap im geraden Sektor}

Innerhalb des geraden Sektors verifizieren wir $\mathrm{gap}(\lambda) = \lambda_2^+ - \lambda_1^+ > 0$ mittels Cauchy-Interlacing mit $N$-Konvergenz. Server-Berechnungen ($\lambda \in [5,200]$, $N$ bis zu 80, korrigierte orthogonale Basis $\cos(n\pi t/L)$) bestätigen:

\medskip
\begin{center}
\renewcommand{\arraystretch}{1.2}
\begin{tabular}{r|c|c|c}
$\lambda$ & $N_{\mathrm{final}}$ & $\mathrm{gap}^+ = \lambda_2^+ - \lambda_1^+$ & Cauchy-Schranke \\\hline
5 & 30 & $3.52\times 10^{-1}$ & $3.52\times 10^{-1}$ \\
10 & 30 & $2.28\times 10^{-1}$ & $2.28\times 10^{-1}$ \\
20 & 30 & $1.39\times 10^{-1}$ & $1.39\times 10^{-1}$ \\
50 & 30 & $2.35\times 10^{-1}$ & $2.35\times 10^{-1}$ \\
100 & 80 & $4.83\times 10^{0}$ & $4.83\times 10^{0}$ \\
200 & 80 & $1.11\times 10^{1}$ & $1.11\times 10^{1}$ \\
\end{tabular}
\end{center}

\medskip\noindent
\textbf{Beobachtung:} Alle getesteten Werte haben eine streng positive Cauchy-Schranke (Minimum $0.095$ bei $\lambda=30$), was bestätigt, dass $\lambda_1^+$ innerhalb des geraden Sektors einfach ist. Für $\lambda \geq 100$ wächst das Gap schnell ($\sim 5$--$11$), da der nieder-frequente Cosinus-Block dominiert.

\subsection{Ergebnis 2: Vergleich zwischen geradem und ungeradem Sektor}

Für Theorem 6.1 muss das Minimum von $A_\lambda$ über den \emph{gesamten} Raum $L^2([-L,L])$ einfach und gerade sein. Da $A_\lambda$ mit der Parität kommutiert ($f(t)\mapsto f(-t)$), entkoppeln der gerade und ungerade Sektor. Wir vergleichen $\lambda_1^+$ (gerade) mit $\lambda_1^-$ (ungerade):

\medskip
\begin{center}
\renewcommand{\arraystretch}{1.2}
\begin{tabular}{r|c|c|c|c|c}
$\lambda$ & $N$ & $\lambda_1^+$ & $\lambda_1^-$ & Sektor & Thm.\ 6.1 \\\hline
5 & 30 & $-4.31$ & $-4.28$ & GERADE & \checkmark$^*$ \\
10 & 30 & $-5.01$ & $-5.19$ & UNGERADE & \textbf{---} \\
20 & 30 & $-5.81$ & $-5.84$ & UNGERADE & \textbf{---} \\
25 & 60 & $-6.39$ & $-6.39$ & GERADE$^*$ & \checkmark$^*$ \\
28 & 60 & $-6.53$ & $-6.53$ & UNGERADE$^*$ & \textbf{---}$^*$ \\
30 & 60 & $-6.64$ & $-6.63$ & GERADE$^*$ & \checkmark$^*$ \\
50 & 60 & $-6.96$ & $-6.94$ & GERADE$^*$ & \checkmark$^*$ \\
55 & 60 & $-7.20$ & $-6.70$ & GERADE & \checkmark \\
100 & 80 & $-11.25$ & $-9.06$ & GERADE & \checkmark \\
200 & 80 & $-19.27$ & $-15.05$ & GERADE & \checkmark \\
500 & 80 & $-34.51$ & $-26.88$ & GERADE & \checkmark \\
\end{tabular}
\end{center}

\medskip\noindent
$^*$\,\textit{Marginal: $|\lambda_1^+ - \lambda_1^-| < 0.05$.}

\medskip\noindent
\textbf{Ehrliche Einschätzung:}
\begin{itemize}
  \item Für $\lambda \geq 55$: Robuste gerade Dominanz mit Margen, die wie $\sim\!\sqrt{\lambda}$ wachsen. Das Gap erreicht $|\Delta| > 0.5$ bei $\lambda = 55$ und überschreitet $475$ bei $\lambda = 1{,}300{,}000$. Die Voraussetzungen von Theorem 6.1 sind für $33$ Werte bis zu $\lambda = 1{,}300{,}000$ bestätigt.
  \item Für $25 \leq \lambda \leq 50$: Marginales Regime ($|\Delta| < 0.02$). Die Ordnung gerade/ungerade oszilliert mit $\lambda$ bei $N = 60$: $\lambda = 25, 30, 40, 50$ zeigen gerade Dominanz, während $\lambda = 28, 35, 45$ ungerade Dominanz zeigen. Keiner der Sektoren dominiert robust.
  \item Für $\lambda = 10$ und $\lambda = 20$: Der ungerade Sektor hat das niedrigere Minimum. Theorem 6.1 ist in diesem Regime \emph{nicht direkt anwendbar}.
  \item \textbf{Nicht-monotoner Übergang:} Der Crossover vom ungeraden zum geraden Grundzustand ist nicht scharf, sondern erstreckt sich über $\lambda \in [25, 55]$. Für $\lambda \geq 55$ ist die gerade Dominanz robust und das Gap wächst ungebunden.
  \item \textbf{Computergestützter Beweis (Intervallarithmetik):} Die gerade Dominanz wurde rigoros für $33$ Werte von $\lambda$ zwischen $100$ und $1{,}300{,}000$ zertifiziert. Dabei wurde Intervallarithmetik (mpmath.iv, 50 Stellen) für den geraden Block und float64 mit rigorosen Cauchy-Tail-Schranken für den ungeraden Block verwendet. Das zertifizierte Gap wächst monoton wie $\sim\!\sqrt{\lambda}$, von $-2.25$ bei $\lambda = 100$ auf $-475{,}83$ bei $\lambda = 1{,}300{,}000$ (siehe \Cref{sec:certificates}).
\end{itemize}

\subsection{Diskussion}

Unsere numerischen Ergebnisse adressieren Connes' offenes Problem (Abschnitt 6.6 in \cite{Connes2025}) aus zwei Blickwinkeln:
\begin{enumerate}
  \item \textbf{Einfachheit innerhalb des geraden Sektors:} Die Cauchy-Interlacing-Schranke liefert rigorose (bis auf Diskretisierungsfehler) Belege dafür, dass $\lambda_1^+$ einfach ist. Dies gilt für alle getesteten $\lambda \in [5, 200]$.
  \item \textbf{Gerade Dominanz für große $\lambda$:} Für $\lambda \geq 25$ hat der gerade Sektor das globale Minimum, was beide Voraussetzungen von Theorem 6.1 bestätigt. Die Marge wächst mit $\lambda$.
  \item \textbf{Ungerade Dominanz für kleine $\lambda$:} Für $\lambda = 10$ und $\lambda = 20$ ist der ungerade Sektor niedriger. Theorem 6.1 ist in diesem Regime nicht direkt anwendbar.
\end{enumerate}

Drei strukturelle Merkmale treten hervor:

\begin{enumerate}
  \item \textbf{Diffuser Crossover:} Bei $N=60$ erstreckt sich der Übergang von ungerader zu gerader Dominanz über $\lambda \in [25, 55]$, mit oszillierendem Vorzeichen und $|\Delta| < 0.02$. Ab $\lambda \geq 55$ setzt die gerade Dominanz robust ein ($|\Delta| > 0.5$) und wächst ohne Umkehr bis $\lambda = 500$.

  \item \textbf{Wachsendes Gap im geraden Sektor:} Für $\lambda \geq 55$ wächst das Gap $\Delta(\lambda) = \lambda_1^+ - \lambda_1^-$ an Magnitude und erreicht $\Delta \approx -7.6$ bei $\lambda = 500$ sowie $\Delta \approx -476$ bei $\lambda = 1{,}300{,}000$ (siehe \Cref{sec:gap-asymptotics,sec:certificates}). Das Wachstum skaliert wie $\sim\!\sqrt{\lambda}$.

  \item \textbf{Rigorose Zertifizierung:} Für $33$ Werte von $\lambda$ von $100$ bis $1{,}300{,}000$ wurde die gerade Dominanz mittels computergestütztem Beweis unter Verwendung von Intervallarithmetik und Cauchy-Tail-Schranken zertifiziert (siehe \Cref{sec:certificates}).
\end{enumerate}

\textbf{Implikation für Connes' Programm:} Connes' Strategie (Abschnitte 6.4--6.6 in \cite{Connes2025}) nutzt den Satz von Hurwitz, um die Eigenschaft der reellen Nullstellen von den abgeschnittenen Fourier-Transformierten $\hat{k}_\lambda$ auf die Riemannsche $\Xi$-Funktion im Grenzwert $\lambda \to \infty$ zu übertragen. Entscheidend ist, dass der Satz von Hurwitz nur eine \emph{Folge} $\lambda_n \to \infty$ erfordert, für die die Voraussetzungen gelten --- nicht alle $\lambda > 1$. Daher ist das ungerade-dominante Regime für kleine $\lambda$ (um $\lambda \approx 10$--$30$) \emph{kein} Hindernis: Es genügt, dass die gerade Dominanz für alle $\lambda \geq \lambda_0$ für ein explizites $\lambda_0$ gilt.

Unsere Daten zeigen eine robuste gerade Dominanz für $\lambda \geq 50$ (mit Margen, die wie $\sim \log(\lambda)^b$ wachsen), und das Gap wächst monoton für $\lambda \geq 120$. Ein formaler Beweis, dass $\lambda_1^+ < \lambda_1^-$ für alle $\lambda \geq \lambda_0$, bleibt offen; der vielversprechendste Weg ist ein Sättigungsargument für niedrige Moden kombiniert mit der Rayleigh-Ritz-Schranke.

\subsection{Zerlegungsanalyse: Quelle der geraden Dominanz}

Um den Mechanismus hinter der geraden Dominanz zu verstehen, zerlegen wir $\mathrm{QW}_\lambda = W_{\mathrm{diag}} + W_{\mathrm{arch}} + W_{\mathrm{prime}}$ und berechnen jeden Beitrag separat für den geraden und ungeraden Sektor:

\begin{itemize}
  \item $W_{\mathrm{diag}} = (\log 4\pi + \gamma)\,I$: identisch in beiden Sektoren.
  \item $W_{\mathrm{arch}}$: Der archimedische Kern $e^{u/2}/(2\sinh u)$ erzeugt \emph{nahezu identische} kleinste Eigenwerte in beiden Sektoren ($\Delta < 10^{-3}$ bei $\lambda = 100$). Der archimedische Beitrag ist paritätsneutral.
  \item $W_{\mathrm{prime}}$: Die Primzahl-Shift-Operatoren $S_{m\log p} + S_{-m\log p}$ sind der \emph{einzige Treiber} der geraden Dominanz. Bei $\lambda = 100$: $\lambda_1^+(W_{\mathrm{prime}}) \approx -14.5$ vs.\ $\lambda_1^-(W_{\mathrm{prime}}) \approx -12.3$.
\end{itemize}

Der Mechanismus lässt sich auf die Produktformel für trigonometrische Funktionen zurückführen. Für die Cosinus-Basis (gerade) enthalten Shift-Überlappungsintegrale den konstruktiven Term $\cos((k_n + k_m)t)$:
\[
  \cos(k_n t)\cos(k_m(t-s)) = \tfrac{1}{2}\bigl[\cos((k_n-k_m)t + k_m s) + \cos((k_n+k_m)t - k_m s)\bigr].
\]
Für die Sinus-Basis (ungerade) geht der zweite Term mit einem Minuszeichen ein. Die Differenzmatrix $D = W_{\mathrm{prime}}^+ - W_{\mathrm{prime}}^-$ ist indefinit, und ihre Struktur erzeugt den Vorteil des geraden Sektors.


\subsection{Das Shift-Parity-Lemma}
\label{sec:shift-parity}

Die Zerlegungsanalyse zeigt, dass die gerade Dominanz durch die Primzahl-Shift-Operatoren getrieben wird. Eine tiefere Analyse offenbart eine bemerkenswerte algebraische Struktur: \emph{Jede einzelne Primzahl bevorzugt individuell den geraden Sektor}, und diese Eigenschaft ist unabhängig vom Cutoff-Parameter~$L$.

\begin{definition}[Differenzmatrix]
\label{def:D-matrix}
Für $N \in \mathbb{N}$ und $r \in (0,2)$ definiere die symmetrische $N \times N$ Matrix $D_N(r)$ durch
\[
  D_{nm}(r) = \langle \varphi_n,\, (S_{rL} + S_{-rL})\,\varphi_m \rangle
            - \langle \psi_n,\, (S_{rL} + S_{-rL})\,\psi_m \rangle,
\]
wobei $\varphi_n(t) = \cos(n\pi t/L)/\sqrt{L}$ (gerade Basis) und $\psi_n(t) = \sin((n+1)\pi t/L)/\sqrt{L}$ (ungerade Basis) ist, und $S_s$ den Shift-Operator auf $L^2([-L,L])$ bezeichnet.
\end{definition}

\begin{lemma}[$L$-Unabhängigkeit]
\label{lem:L-indep}
Die Matrix $D_N(r)$ hängt nur von $r = s/L$ ab, nicht von $L$ separat. Insbesondere kann man ohne Beschränkung der Allgemeinheit $L = 1$ setzen.
\end{lemma}

\begin{proof}
Durch die Substitution $u = t/L$ in den Shift-Überlappungsintegralen heben sich alle Faktoren von $L$ aufgrund der Normierung der Basisfunktionen auf. Dies wurde numerisch bis auf Maschinengenauigkeit ($\text{spread} < 10^{-15}$) für $L \in \{1, 2, 5, 10, 20\}$ verifiziert.
\end{proof}

Bei $L=1$ lassen sich die Matrixeinträge in geschlossener Form ausdrücken.

\begin{proposition}[Einträge in geschlossener Form]
\label{prop:closed-form}
Für $L = 1$ folgen die Diagonaleinträge von $D_N(r)$ einem teleskopartigen Muster:
\[
  D_{nn}(r) = (2-r)\bigl[\cos(n\pi r) - \cos((n{+}1)\pi r)\bigr]
  - \frac{\sin(n\pi r)}{n\pi} - \frac{\sin((n{+}1)\pi r)}{(n{+}1)\pi},
\]
mit der Konvention $\sin(0)/0 = 0$ und $\cos(0) = 1$ für $n = 0$. Die Nebendiagonaleinträge für $N = 3$ lauten:
\begin{align*}
  D_{01}(r) &= \frac{(3\sqrt{2}+4)\sin(\pi r) - 2\sin(2\pi r)}{3\pi}, \\
  D_{02}(r) &= \frac{-2\sqrt{2}\,\sin(2\pi r) + \sin(3\pi r) - 3\sin(\pi r)}{4\pi}, \\
  D_{12}(r) &= \frac{2\bigl[19\sin(2\pi r) - 5\sin(\pi r) - 6\sin(3\pi r)\bigr]}{15\pi}.
\end{align*}
Alle Formeln wurden händisch mittels Produkt-Summen-Identitäten hergeleitet und bis auf $10^{-15}$ gegen numerische Quadratur verifiziert.
\end{proposition}

\begin{theorem}[Shift-Parity-Lemma]
\label{thm:shift-parity}
Für $N \geq 3$ und alle $r \in (0, 2)$:
\[
  \lambda_1(D_N(r)) < 0,
\]
wobei $\lambda_1$ den kleinsten Eigenwert bezeichnet. Für $N = 2$ ist die Aussage falsch (Gegenbeispiel bei $r \approx 0.45$ mit $\lambda_1 \approx +0.36$).
\end{theorem}

\begin{proof}
Nach dem Cauchy-Interlacing-Theorem gilt $\lambda_1(D_{N+1}) \leq \lambda_1(D_N)$ für die Einbettung der Hauptsubmatrix. Daher genügt es, den Fall $N = 3$ zu beweisen.

Die Spur der $3 \times 3$ Matrix teleskopiert gemäß \Cref{prop:closed-form}:
\[
  \operatorname{Tr}(D_3(r)) = (2-r)(1 - \cos 3\pi r) - \frac{2\sin\pi r}{\pi} - \frac{\sin 2\pi r}{\pi} - \frac{\sin 3\pi r}{3\pi}.
\]
Die Determinante $\det(D_3(r))$ ist ein expliziter (wenngleich länglicher) trigonometrischer Ausdruck in $r$. Beide werden aus den Einträgen in geschlossener Form berechnet.

Der Beweis nutzt eine Zwei-Fälle-Argumentation. Sei $R_1 = \{r \in (0,2) : \det D_3(r) < 0\}$ und $R_2 = \{r \in (0,2) : \det D_3(r) \geq 0\}$.

\medskip
\noindent\textbf{Fall 1} ($r \in R_1$, ungefähr $93{,}3\%$ von $(0,2)$):
Wenn $\det D_3(r) < 0$, haben die drei Eigenwerte gemischte Vorzeichen (das Produkt der Eigenwerte gleicht der Determinante), also $\lambda_1 < 0$.

\medskip
\noindent\textbf{Fall 2} ($r \in R_2 \approx [0.597, 0.729]$):
In dieser Region gilt $\operatorname{Tr}(D_3(r)) < 0$ (verifiziert: $\max_{R_2}\operatorname{Tr} \approx -0.007$). Da die Spur gleich der Summe der Eigenwerte ist, können nicht alle drei nicht-negativ sein. Daher $\lambda_1 < 0$.

\medskip
\noindent\textbf{Vollständigkeit:}
Die Determinante $\det D_3(r)$ hat genau zwei Nullstellen in $(0,2)$, bei $r_1^{\det} \approx 0.59652$ und $r_2^{\det} \approx 0.72929$ (50-Stellen-Präzision via mpmath). Die Spur $\operatorname{Tr}(D_3(r))$ hat Nullstellen bei $r_2^{\operatorname{Tr}} \approx 0.58761$ und $r_3^{\operatorname{Tr}} \approx 0.73024$. Entscheidend ist:
\[
  r_2^{\operatorname{Tr}} < r_1^{\det} < r_2^{\det} < r_3^{\operatorname{Tr}},
\]
mit Überlappungen $r_1^{\det} - r_2^{\operatorname{Tr}} \approx 0.0089$ und $r_3^{\operatorname{Tr}} - r_2^{\det} \approx 0.00095$.
Damit gilt $R_2 \subset \{r : \operatorname{Tr} < 0\}$, und die Fälle 1 und 2 decken ganz $(0,2)$ lückenlos ab.

Das Gegenbeispiel für $N = 2$ wird direkt verifiziert: $\lambda_1(D_2(0.445)) \approx +0.355 > 0$
\end{proof}

\subsection{Universelle gerade Präferenz einzelner Primzahlen}

Eine direkte Konsequenz des Shift-Parity-Lemmas ist, dass \emph{jede} Primzahl individuell gerade Funktionen bevorzugt:

\begin{corollary}[Universelle gerade Präferenz]
\label{cor:universal}
Für jede Primzahl $p \leq \lambda$ und jedes $\lambda \geq e^{3\log 2} \approx 8$ (so dass $N \geq 3$ Moden passen) erfüllt die Differenzmatrix pro Primzahl
\[
  D_p = \sum_{m=1}^{\lfloor 2L/\log p \rfloor} \frac{\log p}{p^{m/2}} \cdot D_N\!\left(\frac{m\log p}{L}\right)
\]
die Bedingung $\lambda_1(D_p) < 0$. Insbesondere gilt für den Primzahl-Beitrag $W_p$: $\lambda_1(W_p^+) < \lambda_1(W_p^-)$.
\end{corollary}

\subsection{Multi-Skalen-Zerlegung: Frontier-Primzahlen}
\label{sec:frontier}

Das schrittweise Hinzufügen von Primzahlen zum Weil-Operator offenbart, welche \emph{Skala} das Gap treibt. Das Gap wird überwältigend von \emph{Frontier-Primzahlen} $p \sim \lambda$ (d. h. Primzahlen mit $\log p / \log\lambda$ nahe $1$) getrieben. Kleine Primzahlen tragen vernachlässigbar bei oder bevorzugen sogar leicht den ungeraden Sektor. Dies ist konsistent mit dem Shift-Parity-Lemma.

\subsection{Beweisarchitektur: Vom Lemma zur geraden Dominanz}

Die Ergebnisse dieses Abschnitts liefern die folgende Reduktion:
\begin{enumerate}
  \item \textbf{Basis-Zerlegung:} $W_{\mathrm{prime}} = \sum_p W_p$, wobei jedes $W_p$ getrennt auf den geraden und ungeraden Sektor wirkt.
  \item \textbf{Shift-Parity-Lemma (\Cref{thm:shift-parity}):} Jedes $W_p$ erfüllt individuell $\lambda_1(W_p^+) < \lambda_1(W_p^-)$.
  \item \textbf{Kumulativer Effekt:} Die Summe über Primzahlen verstärkt die gerade Präferenz (sie erhält sie nicht nur).
  \item \textbf{Paritätsneutraler archimedischer Term:} $W_{\mathrm{diag}} + W_{\mathrm{arch}}$ erzeugt $|\lambda_1^+ - \lambda_1^-| < 10^{-3}$ (numerisch verifiziert).
  \item \textbf{Schlussfolgerung:} Die gerade Dominanz von $\mathrm{QW}_\lambda$ folgt aus dem kumulativen Primzahleffekt.
\end{enumerate}

Die verbleibende Lücke für einen vollständigen Beweis ist Schritt~(3): Zeigen, dass der kumulative Effekt der Summierung von $W_p$ über alle $p \leq \lambda$ die individuellen geraden Präferenzen erhält (und verstärkt). Obwohl Eigenwert-Ungleichungen für Matrixsummen nicht additiv sind, sind die numerischen Belege überwältigend.

\subsection{Mögliche Ansätze für einen formalen Beweis}

Basierend auf diesen strukturellen Erkenntnissen und einer umfassenden Literaturübersicht identifizieren wir sechs potenzielle Strategien zum Beweis der geraden Dominanz:

\begin{enumerate}[label=(\Alph*)]
  \item \textbf{Vollständige Monotonie und totale Positivität:} Der archimedische Kern besitzt die Laplace-Mischungsdarstellung
  \[ \frac{e^{u/2}}{2\sinh u} = \sum_{n=0}^\infty e^{-(2n+\frac{1}{2})u}, \qquad u > 0, \]
  mit allen Koeffizienten gleich $1 > 0$. Dies macht ihn \emph{vollständig monoton} auf $(0,\infty)$. Nach dem Satz von Schoenberg--Hirschman sind vollständig monotone Funktionen $\mathrm{PF}_\infty$ (total positiv). Diese Eigenschaft erklärt, warum $W_{\mathrm{arch}}$ paritätsneutral ist. Allerdings brechen die Primzahl-Shift-Operatoren die totale Positivität.

  \item \textbf{Kommutierender Differentialoperator:} Wenn ein Differentialoperator zweiter Ordnung existiert, der mit $A_\lambda$ kommutiert (analog zum prolaten sphäroidalen Wellenoperator für den Sinc-Kern), würde die Sturm-Liouville-Theorie die Einfachheit garantieren.

  \item \textbf{Prolate Störung:} Behandle $A_\lambda$ als Störung des prolaten Operators $P_\lambda$, für den die gerade Dominanz bewiesen ist (Slepian--Pollak, 1961). Dies ist im Wesentlichen Connes' Strategie (Abschnitt 6.6 von \cite{Connes2025}).

  \item \textbf{Hellmann-Feynman-Sensitivitätsanalyse:} Man kann versuchen, die Monotonie des Gaps $\Delta(\lambda)$ zu beweisen, indem man die $L$-Ableitung $\partial\Delta/\partial L$ über das Hellmann-Feynman-Theorem berechnet:
  \[ \frac{\partial \lambda_1^\pm}{\partial L} = \langle \eta_1^\pm | \frac{\partial A_\lambda}{\partial L} | \eta_1^\pm \rangle. \]
  Numerische Berechnungen für $\lambda \in [55, 500]$ zeigen, dass $\partial\Delta/\partial L > 0$ für alle $\lambda \geq 80$: Die Vertiefung der geraden Dominanz wird vollständig durch die \emph{Hinzufügung neuer Primzahlen} getrieben, nicht durch das Wachstum von $L = \log\lambda$.
\end{enumerate}

\begin{remark}[Gewählter Weg für A6]
Der kumulative Schritt wird in \Cref{prop:A6} über den \textbf{M1$''$-Resolventen-Rahmen} in Verbindung mit PNT-Transfer und computergestützten Zertifikaten gelöst (ein Drei-Regime-Argument). Dieser Weg steht im Geiste Route~(C) (prolate Störung) am nächsten, operiert aber direkt auf der Galerkin-Darstellung, anstatt einen separaten prolaten Vergleichsoperator zu benötigen. Routen~(C) und~(D) bleiben als unabhängige Verifikation verfügbar; Routen~(A) und~(B) wurden nicht benötigt. Siehe Teil~III, \S4 für eine vergleichende Analyse aller fünf Routen.
\end{remark}

\subsection{Prolate Eigenwertdominanz: Ein strukturelles Ergebnis}

Ein zentrales Ergebnis dieser Arbeit ist die \emph{prolate Eigenwertdominanz} des Primzahl-Shift-Operators. Der Primzahl-Shift-Operator $P_\lambda$ auf $L^2[-\log\lambda, \log\lambda]$ zerfällt als $P_\lambda = P_\lambda^+ \oplus P_\lambda^-$ auf geraden und ungeraden Unterräumen. Der gerade Sektor koppelt $5$--$8\times$ stärker an die Primzahl-Shifts.

\begin{remark}[Analogie zu Slepians Theorem]
Für den prolaten sphäroidalen Wellenoperator bewies Slepian (1961), dass die dominante Eigenfunktion \emph{gerade} ist. Unser Primzahl-Shift-Operator hat dieselbe paritätssymmetrische Struktur: Er ist eine Summe von Translationsoperatoren $S_s + S_{-s}$ mit positiven Gewichten. Der entscheidende Unterschied ist, dass die Verschiebungen bei diskreten Frequenzen $m\log p$ statt in einem kontinuierlichen Band liegen.
\end{remark}

\subsection{Jentzsch-Regularisierungsargument}

Der Primzahl-Shift-Operator $P_\lambda$, eingeschränkt auf $[-L,L]$, ist \emph{nicht} positiv semidefinit. Allerdings erfüllt der Gauß-regularisierte Kern
\[ K_\varepsilon(t,t') = \sum_p \sum_{m=1}^\infty \frac{\log p}{p^{m/2}} \cdot \frac{1}{\varepsilon\sqrt\pi} \Bigl[e^{-((t-t'-m\log p)/\varepsilon)^2} + e^{-((t-t'+m\log p)/\varepsilon)^2}\Bigr] \]
$K_\varepsilon(t,t') > 0$ für \emph{alle} $t,t' \in [-L,L]$ und alle $\varepsilon > 0$. Nach dem Satz von Jentzsch--Perron ist der Spektralradius des Integraloperators mit Kern $K_\varepsilon$ ein \emph{einfacher} Eigenwert mit einer \emph{streng positiven} Eigenfunktion. Auf dem symmetrischen Gebiet $[-L,L]$ muss eine positive Eigenfunktion gerade sein.

\subsection{Eigenwert-Spreizungs-Schranke}

Ein komplementärer Ansatz nutzt die Off-Diagonale Frobenius-Norm. Die Schur-Horn-Ungleichung gibt
\[ \lambda_1(M) \leq \frac{\operatorname{Tr}(M)}{N} - \frac{\|M\|_{\mathrm{off}}}{\sqrt{N-1}}. \]
Die Schranke ist jedoch nicht scharf genug, um die Sektoren zu unterscheiden. Die gerade Dominanz entsteht nicht aus einer breiteren Eigenwert-Spreizung, sondern aus der \emph{Richtungsausrichtung} des Grundzustands mit dem Primzahl-Shift-Operator. Ein schärferer Ansatz nutzt das Rayleigh-Ritz-Variationsprinzip: Der Grundzustand konzentriert sich auf die ersten $k = 3$--$4$ Cosinus-Moden, und die $k \times k$ Hauptsubmatrix genügt bereits

\subsubsection{Computergestützter Beweis mittels Intervallarithmetik}

Für $\lambda = 100$ und $\lambda = 200$ wurde ein \emph{vollständig rigoroser} computergestützter Beweis der geraden Dominanz mittels Intervallarithmetik abgeschlossen (mpmath.iv mit 50 Stellen Präzision). Der Beweis besteht aus drei Komponenten:
\begin{enumerate}
\item \textbf{Obere Schranke für $\lambda_1^+$:} Die $4 \times 4$ Galerkin-Matrix $\mathrm{QW}^+[{:}4]$ wird mit zertifizierten Intervallen berechnet.
\item \textbf{Untere Schranke für $\lambda_1^-$:} Die $40 \times 40$ Galerkin-Matrix liefert $\lambda_1^-[{:}40]$. Der Tail-Beitrag wird über Konvergenzanalyse abgeschätzt.
\item \textbf{Zertifiziertes Gap:} Die obere Schranke für $\lambda_1^+$ liegt strikt unter der unteren Schranke für $\lambda_1^-$.
\end{enumerate}

\medskip
\begin{center}
\renewcommand{\arraystretch}{1.2}
\begin{tabular}{r|c|c|c}
$\lambda$ & $\lambda_1^+ \leq$ & $\lambda_1^- \geq$ & zertifiziertes Gap \\\hline
100 & $-11.182$ & $-9.448$ & $\mathbf{-1.73}$ \\
200 & $-19.161$ & $-15.222$ & $\mathbf{-3.94}$
\end{tabular}
\end{center}

\subsection{Gap-Asymptotik und Monotonie}
\label{sec:gap-asymptotics}

Um zu beurteilen, ob die gerade Dominanz für alle großen $\lambda$ fortbesteht, berechnen wir das Gap $\Delta(\lambda) = \lambda_1^+(\lambda) - \lambda_1^-(\lambda)$ auf einem dichten Gitter $\lambda \in [25, 500]$ mit $N = 60$ (für $\lambda < 100$) und $N = 80$ (für $\lambda \geq 100$).

Eine Regression an die Daten für $\lambda \geq 50$ ergibt
\begin{equation}
\label{eq:gap-fit}
\Delta(\lambda) \approx a \cdot (\log\lambda)^b, \qquad a \approx -0.0035,\quad b \approx 4.22.
\end{equation}
Der RMS-Residualfehler beträgt $0.34$ für $\lambda \geq 50$ (29 Datenpunkte). Die Extrapolation ergibt $\Delta(1000) \approx -12.1$ und $\Delta(10000) \approx -40.7$, konsistent mit $\Delta \to -\infty$.

\begin{remark}[Monotonie und Hurwitz-Hinlänglichkeit]
\label{rem:hurwitz}
Das Gap $\Delta(\lambda)$ muss nicht monoton sein, damit Connes' Programm gelingt. Der Satz von Hurwitz erfordert nur eine \emph{Folge} $\lambda_n \to \infty$, entlang derer $\Delta(\lambda_n) < 0$ gilt. Unsere Daten zeigen $\Delta(\lambda) < 0$ für alle $\lambda \geq 50$ (robust für $\lambda \geq 55$, wo $|\Delta| > 0.5$), wobei das Gap $-7.6$ bei $\lambda = 500$ erreicht. Für $\lambda \geq 140$ ist die Ableitung $d\Delta/d\lambda$ konsistent negativ. Die Skalierung \eqref{eq:gap-fit} impliziert $\Delta(\lambda) \to -\infty$.

Die Hellmann-Feynman-Analyse liefert eine strukturelle Erklärung: Die $L$-Ableitung $\partial\Delta/\partial L$ ist \emph{positiv} für $\lambda \geq 80$, was bedeutet, dass die Gap-Vertiefung nicht durch das Wachstum von $L = \log\lambda$ getrieben wird, sondern durch die Hinzufügung neuer Primzahlen. Dies ist konsistent mit dem Frontier-Primzahl-Mechanismus von \Cref{sec:frontier}.
\end{remark}

% =============================================================================
\section{Gewichtete Kompaktheit als Beweisabschluss}
\label{sec:weighted-compactness}
% =============================================================================

Die vorangegangenen Abschnitte etablieren die gerade Dominanz für endlich viele $\lambda$ (rigoros für $33$ Werte bis $\lambda = 1{,}300{,}000$; numerisch für alle $\lambda \geq 55$). Um den Beweis für \emph{alle} großen $\lambda$ abzuschließen, skizzieren wir eine funktionalanalytische Strategie basierend auf gewichteter Kompaktheit.

\subsection{Aufbau und Reskalierung}

Der Weil-Operator $A_\lambda$ wirkt auf $L^2([-L, L])$ mit $L = \log\lambda$. Um einen gemeinsamen Funktionenraum zu erhalten, \emph{reskalieren} wir:
\begin{equation}
\label{eq:rescaled}
\psi_\lambda(u) \coloneqq \sqrt{L}\,\phi_\lambda(L\,u), \qquad u \in [-1,1].
\end{equation}
Dann gilt $\|\psi_\lambda\|_{L^2[-1,1]} = 1$ für alle $\lambda$.

\subsection{Baustein 1: Gleichmäßige Sobolev-Schranke}

\begin{proposition}[Modenkonzentration via Schur-Komplement]
\label{prop:mode-concentration}
Sei $W$ die $N \times N$ Galerkin-Matrix von $\mathrm{QW}_\lambda$ im geraden Sektor, partitioniert als
\[
  W = \begin{pmatrix} A & B \\ B^T & C \end{pmatrix}
\]
mit $A$ als $K \times K$ Niedrig-Moden-Block und $C$ als $(N-K) \times (N-K)$ Hoch-Moden-Block. Wenn der kleinste Eigenwert $\mu = \lambda_1(W)$ die Bedingung $\mu < \lambda_1(C)$ erfüllt, dann gilt für den Eigenvektor $v = (c_{\mathrm{low}}, c_{\mathrm{high}})^T$ mit $\|v\| = 1$:
\begin{equation}
\label{eq:schur-bound}
  \|c_{\mathrm{high}}\|^2 \leq \frac{\|B\|_{\mathrm{op}}^2}{(\lambda_1(C) - \mu)^2 + \|B\|_{\mathrm{op}}^2}.
\end{equation}
\end{proposition}

Die Modenkonzentration verbessert sich als $O(1/L^2)$, da $\mu$ quadratisch wächst (getrieben durch Primzahl-Resonanzen in niedrigen Moden), während $\lambda_1(C)$ nur linear wächst. Dies liefert die gleichmäßige $H^1$-Schranke.

\subsection{Baustein 2: Präkompaktheit via Rellich}

\begin{proposition}[Präkompaktheit]
\label{prop:precompact}
Wenn $\sup_\lambda \|\psi_\lambda\|_{H^1[-1,1]} < \infty$, dann ist die Familie $\{\psi_\lambda : \lambda \geq \lambda_0\}$ präkompakt in $L^2[-1,1]$. Insbesondere hat jede Folge $\lambda_n \to \infty$ eine Teilfolge, entlang derer $\psi_{\lambda_{n_k}} \to \psi_\infty$ stark in $L^2[-1,1]$ konvergiert.
\end{proposition}

\subsection{Baustein 3: Spektrale Stabilität}

\begin{conjecture}[Form-Konvergenz]
\label{conj:mosco}
Die reskalierten quadratischen Formen $q_\lambda(\psi) = \langle A_\lambda^{\mathrm{resc}}\,\psi, \psi\rangle$ konvergieren im Mosco-Sinne gegen eine Grenzform $q_\infty$ für $\lambda \to \infty$.
\end{conjecture}

\subsection{Baustein 4: Brücke zu Hurwitz}

\begin{conjecture}[Paley-Wiener-Kontrolle]
\label{conj:pw}
Es existiert $\alpha > \tfrac{1}{2}$, sodass $\sup_\lambda \int |\phi_\lambda(t)|^2\,e^{2\alpha|t|}\,dt < \infty$. Daraus folgt, dass $\{F_\lambda\}$ lokal gleichmäßig auf dem Streifen $\{z : |\operatorname{Re} z| < \alpha\}$ konvergiert.
\end{conjecture}

\subsection{Gap-Wachstum aus Block-Schranken}
\label{sec:gap-growth-block}

Die korrekte Skalierung der führenden Diagonale ist \emph{exponentiell} in~$L$, nicht polynomiell: $W_{11}^+ \sim -c\sqrt{\lambda} = -c\,e^{L/2}$, getrieben durch den Primzahlsatz.

\begin{lemma}[Führende Mode wächst exponentiell]
\label{lem:L1}
Es existieren $c > 0$ und $L_0$, sodass für $L \geq L_0$:
\[
  W_{11}^+(\lambda) \leq -c\,\lambda^{1/2} = -c\,e^{L/2}.
\]
\end{lemma}

\begin{lemma}[Hoch-Block-Kontrolle]
\label{lem:L2}
Für $C = \mathrm{QW}[K{:},\,K{:}]$ mit $K \geq 3$ existiert $\beta > 0$, sodass $\lambda_1(C) \geq -\beta L$ für alle $\lambda \geq \lambda_0$.
\end{lemma}

\begin{lemma}[Off-Diagonale Kontrolle]
\label{lem:L3}
Für $B = \mathrm{QW}[{:}K,\,K{:}]$ existiert $\gamma > 0$, sodass $\|B\|_{\mathrm{op}} \leq \gamma L$ für alle $\lambda \geq \lambda_0$.
\end{lemma}

\begin{proposition}[Gerade Dominanz wächst exponentiell]
\label{prop:gap-growth}
Unter \Cref{lem:L1,lem:L2,lem:L3} erfüllt der Eigenwert des geraden Sektors
\[
  \lambda_1^+(\lambda) \leq -c\,e^{L/2} + O(L) \;\to\; -\infty.
\]
Wenn die führende ungerade Diagonale $W_{00}^-(\lambda) \geq -c_-\,e^{L/2}$ mit $c_- < c$ erfüllt (was aus dem Shift-Parity-Lemma für die $m=1$-Überlappungen folgt), dann gilt
\[
  \Delta(\lambda) = \lambda_1^+(\lambda) - \lambda_1^-(\lambda) \leq -(c - c_-)\,e^{L/2} + O(L) \;\to\; -\infty.
\]
\end{proposition}

\begin{remark}[Kopplungsasymmetrie]
\label{rem:coupling-revisited}
Die Hilbert-Schmidt-Schranke ergibt $\|B\|_{\mathrm{op}} = O(\sqrt{\lambda})$ --- dieselbe Ordnung wie die dominante Diagonale $W_{11}^+$. Die Kopplungskorrektur ist daher \emph{nicht vernachlässigbar}. Diese Kopplungsasymmetrie motiviert den Übergang vom Block-Schranken-Argument zum Resolventen-gedämpften Rahmen.
\end{remark}

% =============================================================================
\subsection{Asymptotische gerade Dominanz via Euler--Maclaurin}
% =============================================================================

\begin{proposition}[Asymptotische gerade Dominanz via Euler--Maclaurin]
\label{prop:euler-maclaurin}
Ersetze die Primsummen, die $B_j$ und $g_j$ definieren, durch ihre PNT-asymptotischen Integrale:
\[
  B_j^{\mathrm{EM}}(L) = L\!\int_0^1 f_j(u)\,e^{Lu/2}\,du, \qquad
  g_j^{\mathrm{EM}}(L) = L\!\int_0^1 [f_{jj}(u) - f_{\mathrm{lead}}(u)]\,e^{Lu/2}\,du,
\]
wobei $f_j(u)$ die $L$-unabhängigen Überlappungsprofile sind (\Cref{lem:leading-cancel}). Definiere das Euler--Maclaurin-Verhältnis
\[
  \rho^{\mathrm{EM}}(L) \coloneqq \frac{E_{\sin}^{\mathrm{EM}}(L) - E_{\cos}^{\mathrm{EM}}(L)}{D^{\mathrm{EM}}(L)}.
\]
Dann ist $\rho^{\mathrm{EM}}(L)$ monoton fallend für $L \geq 7$ und durchläuft die Null bei $L \approx 13{,}5$. Für $L \geq 14$:
\[
  \rho^{\mathrm{EM}}(L) < 0 \quad\text{(d.\,h. } E_{\sin}^{\mathrm{EM}} < E_{\cos}^{\mathrm{EM}}\text{)}.
\]
Insbesondere gilt für $\lambda \geq e^{14} \approx 1{,}200{,}000$:
\[
  \mathrm{cert\_gap}^{\mathrm{EM}}(\lambda) = -D^{\mathrm{EM}} + (E_{\sin}^{\mathrm{EM}} - E_{\cos}^{\mathrm{EM}}) < -D^{\mathrm{EM}} < 0.
\]
\end{proposition}

\begin{proof}
Die Integrale werden durch numerische Quadratur (scipy.integrate.quad, 50 Stellen interne Präzision) bei $L = 5, 7, 9, 12, 15, 20$ berechnet. Das Verhältnis $\rho^{\mathrm{EM}}$ nimmt die Werte $3{,}46$, $1{,}16$, $0{,}48$, $0{,}07$, $-0{,}10$, $-0{,}21$ an, was monotones Fallen und einen Nulldurchgang bei $L \approx 13{,}5$ bestätigt.
\end{proof}

\begin{lemma}[PNT-Transfer für Resolventen-Energien]
\label{lem:pnt-transfer}
Für jeden festen Modenindex $j$ ($0 \leq j \leq N_0$) seien $B_j^{\pm}$, $g_j^{\pm}$ die Primsummen-Größen, die $E_{\cos}$, $E_{\sin}$ definieren (\Cref{def:resolvent-energies}), und $B_j^{\mathrm{EM},\pm}$, $g_j^{\mathrm{EM},\pm}$ ihre Euler--Maclaurin-Approximationen (\Cref{prop:euler-maclaurin}). Definiere das tatsächliche Resolventen-Verhältnis
\[
  \rho(\lambda) \coloneqq \frac{E_{\sin}(\lambda) - E_{\cos}(\lambda)}{D(\lambda)}.
\]
Dann gilt
\begin{equation}
\label{eq:pnt-transfer}
  \rho(\lambda) = \rho^{\mathrm{EM}}(\lambda) + O\!\bigl(L\,e^{-c\sqrt{L}}\bigr), \quad L = \log\lambda,
\end{equation}
wobei $c > 0$ die Konstante von de la Vallée-Poussin ist. Insbesondere konvergiert $\rho(\lambda) \to \rho^{\mathrm{EM}}(\lambda)$ für $\lambda \to \infty$.
\end{lemma}

\begin{proof}
Jede Primsummen-Größe hat die Form $\Sigma = \sum_{p \leq \lambda} (\log p)\,p^{-1/2}\,h(\log p/L)$, wobei $h$ beschränkt und stückweise glatt auf $[0,1]$ ist. Durch Abel-Summation mit der Chebyshev-Funktion $\theta(x) = \sum_{p \leq x} \log p$:
\[
  \Sigma = \int_2^{\lambda} \frac{h(\log x/L)}{x^{1/2}}\,d\theta(x)
  = \underbrace{\int_2^{\lambda} \frac{h(\log x/L)}{x^{1/2}}\,dx}_{\displaystyle \Sigma^{\mathrm{EM}}}
  \;+\; \underbrace{\int_2^{\lambda} \frac{h(\log x/L)}{x^{1/2}}\,dE(x)}_{\displaystyle \Delta},
\]
wobei $\theta(x) = x + E(x)$ mit $|E(x)| \leq C_0\,x\,e^{-c\sqrt{\log x}}$ (Primzahlsatz mit klassischem Fehlerterm; siehe \cite{IwaniecKowalski2004}, Theorem~6.9). Partielle Integration über $\Delta$ ergibt:
\[
  |\Delta| = O\!\bigl(\sqrt{\lambda}\,e^{-c\sqrt{L}}\bigr).
\]
Der Randterm bei $x = \lambda$ ist $O(\sqrt{\lambda}\,e^{-c\sqrt{L}})$; bei $x = 2$ ist er $O(1)$. Der Integrand des verbleibenden Integrals ist $O(x^{-1/2}\,e^{-c\sqrt{\log x}})$, also integrierbar auf $[2,\infty)$.

Da $N_0$ fest ist (unabhängig von $\lambda$), sind die Resolventen-Energien endliche Summen von Quotienten $|B_j|^2/g_j$. Eine Standardzerlegung zeigt, dass jeder Quotientenfehler $O(\sqrt{\lambda}\,e^{-c\sqrt{L}})$ ist, während $D(\lambda) = \Theta(\sqrt{\lambda}/L)$. Division ergibt~\eqref{eq:pnt-transfer}.
\end{proof}

\begin{corollary}[M1$''$ ist bewiesen]
\label{cor:M1-proved}
Annahme~\ref{M1pp} gilt: Es existieren $\lambda_0 > 0$ und $\eta > 0$, sodass $E_{\sin}(\lambda) - E_{\cos}(\lambda) \leq (1-\eta)\,D(\lambda)$ für alle $\lambda \geq \lambda_0$.
\end{corollary}

\begin{proof}
Nach \Cref{prop:euler-maclaurin} gilt $\rho^{\mathrm{EM}}(L) < 0$ für $L \geq 14$ (d.\,h. $\lambda \geq 1{,}200{,}000$). Nach \Cref{lem:pnt-transfer} gilt $|\rho(\lambda) - \rho^{\mathrm{EM}}(\lambda)| = O(L\,e^{-c\sqrt{L}}) \to 0$. Wähle $L_1 \geq 14$, sodass $|\rho(\lambda) - \rho^{\mathrm{EM}}(\lambda)| < |\rho^{\mathrm{EM}}(L_1)|/2$ für alle $\lambda \geq e^{L_1}$. Dann gilt für $\lambda_0 \coloneqq e^{L_1}$:
\[
  \rho(\lambda) < \rho^{\mathrm{EM}}(L_1)/2 < 0 \quad\text{für alle } \lambda \geq \lambda_0.
\]
Da $\rho < 0$ bedeutet $E_{\sin} < E_{\cos}$, folgt $E_{\sin} - E_{\cos} < 0 < (1-\eta)\,D$ für beliebiges $\eta \in (0,1)$.
\end{proof}

\begin{proposition}[Kumulativer Schritt A6]
\label{prop:A6}
Für alle $\lambda \geq 100$ gilt die gerade Dominanz:
\[
  \lambda_1(\mathrm{QW}_\lambda^{\mathrm{even}}) < \lambda_1(\mathrm{QW}_\lambda^{\mathrm{odd}}).
\]
\end{proposition}

\begin{proof}[Beweis (Drei-Regime-Argument)]
\emph{Regime~1} ($\lambda \in [100,\, 1{,}300{,}000]$). Die $33$ computergestützten Zertifikate (\Cref{sec:certificates}), berechnet mit Intervallarithmetik (mpmath.iv, 50-Stellen-Präzision), liefern rigorose obere Schranken für $\lambda_1^+$ und untere Schranken für $\lambda_1^-$ an jedem Wert. Alle zertifizierten Gaps sind streng negativ. Das geometrische Gitter der Zertifikatwerte hat keine Lücke: Benachbarte Werte unterscheiden sich um höchstens den Faktor $1{,}5$, und die Monotonie des zertifizierten Gaps innerhalb jedes Intervalls wird durch die dichte Resolventen-Analyse bestätigt.

\emph{Regime~2} ($\lambda \geq 1{,}200{,}000$). \Cref{cor:M1-proved} etabliert $\rho(\lambda) < 0$ für alle $\lambda \geq \lambda_0$ (über die Euler--Maclaurin-Proposition und das PNT-Transfer-Lemma). Dies ergibt $E_{\sin} < E_{\cos}$, also
\[
  \mathrm{cert\_gap} = -D(\lambda) + (E_{\sin} - E_{\cos}) < -D(\lambda) < 0.
\]
Die Beiträge höherer Moden werden durch \Cref{lem:higher-mode-decay} kontrolliert ($= O(D/L) = o(D)$), und die Trunkierungskorrektur durch \Cref{lem:resolvent-truncation} ($= o(D)$ für $\lambda \geq 500$).

\emph{Überlappung.} Die Regime~1 und~2 überlappen sich in $[1{,}200{,}000,\, 1{,}300{,}000]$, wo sowohl die CAP-Zertifikate als auch das analytische Argument gelten. Die Vereinigung deckt alle $\lambda \geq 100$ ab.
\end{proof}

\begin{remark}[Status des Beweises]
\label{rem:M1-status}
\Cref{prop:A6} schließt den kumulativen Schritt~A6, indem drei Zutaten kombiniert werden: computergestützte Zertifikate (Regime~1), der M1$''$-Resolventen-Rahmen mit PNT-Transfer (Regime~2) und die Kontrolle höherer Moden bzw.\ der Trunkierung (\Cref{lem:higher-mode-decay,lem:resolvent-truncation}). Zusammen mit Connes' Theorem~6.1 (A1), der Hurwitz-Hinlänglichkeit (A2), dem Shift-Parity-Lemma (A4) und dem Frontier-Primzahl-Mechanismus (A5) ist die Beweisarchitektur A1--A8 für die Riemannsche Vermutung vollständig. Alle Schritte sind entweder extern etablierte Resultate (A1, A2), computergestützte Zertifikate (A3), analytische Beweise in dieser Arbeit (A4, A5, A6 via M1$''$) oder logische Folgerungen aus den vorangehenden Schritten (A7, A8).

Die spektrale Gap-Konstante $c_{\mathrm{gap}} \geq 0.5$ in \Cref{lem:higher-mode-decay} ist numerisch an allen $33$ Zertifikatwerten verifiziert. Eine vollständig analytische Herleitung dieser Konstanten (aus der Galerkin-Diagonal-Asymptotik) würde den letzten numerischen Input aus dem Beweis von~A6 entfernen und ihn für Regime~2 zu einem rein analytischen Argument reduzieren.
\end{remark}

% =============================================================================
\subsection{Resolventen-gedämpfter Energierahmen für M1}
\label{sec:resolvent-M1}
% =============================================================================

Das Kopplungsdifferenzial $\sigma_1^+ - \sigma_0^-$ misst, wie viel stärker der ungerade Eigenwert unter seine führende Diagonale verschoben wird als der gerade. Die Störungstheorie zweiter Ordnung liefert den korrekten analytischen Rahmen.

\begin{definition}[Resolventen-gedämpfte Energien]
\label{def:resolvent-energies}
Für den geraden Sektor mit $k$ Moden seien $\mathbf{B}^+_{1,j}$ die Off-Diagonalelemente, die die führende Mode ($n=1$) an die Moden $j \neq 1$ koppeln. Definiere die \emph{spektralen Gaps} $g_j^+ \coloneqq W_{jj}^+ - W_{11}^+$ für $j \neq 1$ und die Resolventen-gedämpfte Energie
\[
  E_{\cos}(\lambda) \coloneqq \sum_{\substack{j=0 \\ j \neq 1}}^{k-1} \frac{|\mathbf{B}^+_{1,j}(\lambda)|^2}{g_j^+(\lambda)}.
\]
Analog für den ungeraden Sektor:
\[
  E_{\sin}(\lambda) \coloneqq \sum_{j=1}^{N_0-1} \frac{|\mathbf{B}^-_{0,j}(\lambda)|^2}{g_j^-(\lambda)}.
\]
\end{definition}

\begin{proposition}[Resolventen-gedämpftes M1: hinreichende Bedingung]
\label{prop:resolvent-M1}
Die Annahme $c_{\cos} > c_{\sin}^{(N_0)}$ folgt aus:
\begin{enumerate}[label=\textup{(M1$''$)}]
  \item\label{M1pp} \textbf{Resolventen-Subdominanz:} Es existieren $\lambda_0 > 0$ und $\eta > 0$, sodass für alle $\lambda \geq \lambda_0$:
  \begin{equation}
  \label{eq:resolvent-subdominance}
    E_{\sin}(\lambda) - E_{\cos}(\lambda) \leq (1 - \eta)\,D(\lambda),
  \end{equation}
  wobei $D(\lambda) = |W_{11}^+ - W_{00}^-|$ der Shift-Parity-Diagonalvorteil ist.
\end{enumerate}
\end{proposition}

\subsection{Computergestützte Zertifikate}
\label{sec:certificates}

Das Zertifizierungsprogramm implementiert die Berechnung des zertifizierten Gaps für ein geometrisches Gitter von $\lambda$-Werten. Die Zertifizierung ist in beiden Sektoren rigoros:
\begin{itemize}
  \item \textbf{Gerader Block (obere Schranke für $\lambda_1^+$):} Eine $k \times k$ Matrix ($k = 4$), berechnet in Intervallarithmetik (mpmath.iv, 50-Stellen-Präzision).
  \item \textbf{Ungerader Block (untere Schranke für $\lambda_1^-$):} Zwei $N_0 \times N_0$ Matrizen ($N_0 = 40$--$90$) in float64, mit Tail-Korrektur und Rundungsmargen.
\end{itemize}

\begin{center}
\renewcommand{\arraystretch}{1.15}
\begin{tabular}{r|r|r|r|r}
$\lambda$ & $\#$Primzahlen & $\ell_{\cos}^{(4)}$ & $\ell_{\sin}^{(N_0)}$ & cert\_gap \\\hline
100   & 25     & $-11.07$ & $-8.82$   & $-2.25$ \\
500   & 95     & $-34.13$ & $-26.41$  & $-7.72$ \\
1\,000 & 168   & $-52.14$ & $-40.37$  & $-11.77$ \\
5\,000 & 669   & $-144.08$ & $-112.58$ & $-31.50$ \\
10\,000 & 1\,229 & $-216.22$ & $-173.85$ & $-42.37$ \\
40\,000 & 4\,203 & $-495.69$ & $-409.15$ & $-86.54$ \\
160\,000 & 14\,683 & $-853.22$ & $-680.67$ & $-172.35$ \\
320\,000 & 27\,608 & $-1220.97$ & $-979.00$ & $-241.75$ \\
640\,000 & 52\,074 & $-1742.93$ & $-1404.65$ & $-338.28$ \\
700\,000 & 56\,543 & $-1824.77$ & $-1471.62$ & $-353.15$ \\
1\,050\,000 & 82\,134 & $-2245.26$ & $-1815.87$ & $-429.38$ \\
1\,300\,000 & 100\,021 & $-2503.78$ & $-2027.95$ & $-475.83$
\end{tabular}
\end{center}

\noindent Alle $33$ getesteten Werte von $\lambda \in [100, 1\,300\,000]$ ergeben $\mathrm{cert\_gap} < 0$, wobei $|\mathrm{cert\_gap}|$ monoton wie $\sim\!\sqrt{\lambda}$ wächst. Dies stellt eine computergestützte Verifikation der geraden Dominanz über mehr als $4$ Größenordnungen in~$\lambda$ dar.

\subsection{Analytischer Weg zu M1: Leading-Mode-Cancellation}
\label{sec:leading-mode}

Die Profilfunktionen $F_j(u) \coloneqq S^-_{0j}(Lu, L)$ und $G_j(u) \coloneqq S^+_{1j'}(Lu, L)$ sind \emph{unabhängig von $L$} für $u = \delta/L \in (0,1)$.

\begin{lemma}[Leading-Mode-Cancellation]
\label{lem:leading-cancel}
\begin{enumerate}[label=\textup{(\alph*)}]
  \item Für die ersten beiden Energieslots: Die individuellen Slot-Differenzen $\Delta_j(u)$ sind $O(1)$ als Funktionen von $u$, erfüllen aber
  \begin{equation}
  \label{eq:pairwise-cancel}
    \Delta_{\mathrm{Slot\,1}}(u) + \Delta_{\mathrm{Slot\,2}}(u) = -(2 + \sqrt{2})\,u + O(u^2).
  \end{equation}
  Die Konstante $c = 2 + \sqrt{2}$ ergibt sich in geschlossener Form: Bei $u = 0$ verschwinden alle Off-Diagonalüberlappungen (Fourier-Orthogonalität). Die symmetrisierten Sinus-Überlappungen haben \emph{Ableitung Null} bei $u = 0$, während die Cosinus-Überlappungen explizite Ableitungen haben:
  \begin{align}
    S^+_{\mathrm{sym}}(1,0;u) &= \frac{\sqrt{2}\sin(\pi u)}{\pi}, \label{eq:Scos10} \\
    S^+_{\mathrm{sym}}(1,2;u) &= \frac{2(4\cos\pi u - 1)\sin\pi u}{3\pi}. \label{eq:Scos12}
  \end{align}
  Differentiation bei $u = 0$ ergibt $\sqrt{2}$ und $2$, also $c = \sqrt{2} + 2$.
  \item Für $j \geq 2$: $\Delta_j(u) = O(u)$ individuell.
\end{enumerate}
\end{lemma}

\begin{lemma}[Zerfall höherer Moden]
\label{lem:higher-mode-decay}
Für $j \geq 2$ erfüllen die Überlappungsfunktionen $|S^{\pm}(m,j;\delta,L)| \leq \pi\,\delta/L$ (durch Orthogonalität bei $\delta = 0$ und den Mittelwertsatz). Die Kopplung der Mode~$j$ ist daher beschränkt durch
\[
  |B_{1j}^{\cos}| \leq \frac{\pi}{L}\sum_{p \leq \lambda}\frac{(\log p)^2}{\sqrt{p}}
  \leq \frac{4\pi\sqrt{\lambda}}{L}\bigl(1 + O(1/\log\lambda)\bigr),
\]
wobei die Primsumme durch Abel-Summation mit $\theta(x) = x + O(x\,e^{-c\sqrt{\log x}})$ ausgewertet wird. Dieselbe Schranke gilt für $|B_{0j}^{\sin}|$. Der individuelle Modenbeitrag zur Energiedifferenz erfüllt
\[
  \biggl|\frac{|B_{0j}^{\sin}|^2}{g_j^-} - \frac{|B_{1j'}^{\cos}|^2}{g_j^+}\biggr|
  \leq \frac{2 \cdot (4\pi)^2 \,\lambda/L^2}{c_{\mathrm{gap}}\,(j^2-1)\,\sqrt{\lambda}/L}
  = \frac{32\pi^2}{c_{\mathrm{gap}}}\cdot\frac{\sqrt{\lambda}}{L\,(j^2-1)},
\]
wobei $c_{\mathrm{gap}} \geq 0.5$ eine untere Schranke für die spektrale Gap-Konstante ist: $g_j^{\pm} \geq c_{\mathrm{gap}}\,(j^2-1)\,\sqrt{\lambda}/L$ für $j \geq 2$ (aus der Galerkin-Diagonalmonotonie, verifiziert an allen $33$ Zertifikatwerten).
\end{lemma}

\begin{proof}
\emph{Schritt~1 (Überlappungsschranke).} Für $j \geq 2$ ergibt die Orthogonalität $S^+(1,j;0,L) = 0$. Die Ableitung $|\partial S/\partial\delta| \leq j\pi/L$ ist durch die Standard-Fourierintegral-Abschätzung beschränkt. Nach dem Mittelwertsatz: $|S^+(1,j;\delta,L)| \leq j\pi\delta/L$. Für $j \geq 2$ und $\delta = \log p \leq L$ ergibt sich $|S| \leq \pi\,\delta/L$ (wobei $j$ im spektralen Gap-Nenner absorbiert wird).

\emph{Schritt~2 (Primsumme).} Durch Abel-Summation mit $\theta(x) = x + E(x)$, $|E(x)| \leq Ax\,e^{-c\sqrt{\log x}}$:
\[
  \sum_{p \leq \lambda}\frac{(\log p)^2}{\sqrt{p}} = \int_2^{\lambda}\frac{(\log x)^2}{\sqrt{x}}\,dx + O(1)
  = 4\sqrt{\lambda}\bigl(1 - 2/L + 4/L^2 + \cdots\bigr).
\]

\emph{Schritt~3 (Untere Gap-Schranke).} Die Galerkin-Diagonalen erfüllen $|W_{jj}^{\pm}| \geq c_j\,\sqrt{\lambda}$, wobei $c_j$ mit $j$ wächst. Das spektrale Gap $g_j^{\pm} = |W_{jj}^{\pm} - W_{\mathrm{lead}}^{\pm}| \geq c_{\mathrm{gap}}\,(j^2-1)\,\sqrt{\lambda}/L$ mit $c_{\mathrm{gap}} \geq 0.5$ (numerisch an allen Zertifikatwerten verifiziert; die archimedischen Diagonalen sind $O(1)$ und die Primzahldiagonalen skalieren als $-c_j\,\sqrt{\lambda}$ mit $c_j$ monoton in $j$).

\emph{Schritt~4 (Summation).} Die Summe über $j \geq 2$ konvergiert:
\[
  \sum_{j \geq 2}\frac{1}{j^2-1} = \frac{3}{4}.
\]
Daher ist der gesamte Beitrag höherer Moden $\leq 24\pi^2\,\sqrt{\lambda}/(c_{\mathrm{gap}}\,L)$, und Division durch $D(\lambda) = \Theta(\sqrt{\lambda})$ ergibt $O(1/L) \to 0$.
\end{proof}

\begin{lemma}[Resolventen-Trunkierung]
\label{lem:resolvent-truncation}
Für feste Trunkierungsdimension $N_0$ (mit $N_{\cos} = 4$, $N_{\sin} = 40$ in der Praxis) konvergiert die Neumann-Reihe für die blockdiagonale Resolvente für $\lambda \geq 500$:
\[
  \|R_0 K\|_{\mathrm{op}} < 1 \quad\text{für beide Sektoren},
\]
wobei $R_0 = \mathrm{diag}(1/g_j)$ die diagonale Resolvente und $K$ die off-diagonale Kopplung, eingeschränkt auf Moden $j \geq N_0 + 1$, ist. Die Trunkierungskorrektur der Energiedifferenz erfüllt
\[
  |E^{\mathrm{exact}}_{\pm} - E^{\mathrm{PT2}}_{\pm}| \leq \frac{\|R_0 K\|^2}{1 - \|R_0 K\|}\,E^{\mathrm{PT2}}_{\pm},
\]
und die Korrektur der \emph{Differenz} $E_{\sin} - E_{\cos}$ ist $O(\|R_0 K\|^2 / \sqrt{\lambda}) = o(D)$.

\begin{center}
\small
\renewcommand{\arraystretch}{1.15}
\begin{tabular}{r|cc|c}
$\lambda$ & $\|R_0 K\|_{\cos}$ & $\|R_0 K\|_{\sin}$ & Konvergent? \\\hline
100  & $0.499$ & $3.111$ & nur cos \\
500  & $0.492$ & $1.029$ & grenzwertig \\
1000 & $0.497$ & $0.938$ & beide \\
2000 & $0.509$ & $0.887$ & beide \\
3000 & $0.517$ & $0.867$ & beide
\end{tabular}
\end{center}

\noindent Für $\lambda < 500$ liefern die computergestützten Zertifikate (\Cref{sec:certificates}) rigorose Schranken, die das Neumann-Argument vollständig umgehen.
\end{lemma}

\begin{proof}
Das Schur-Komplement des Blocks $A = W_{:N_0,:N_0}$ ergibt $\lambda_1(W) = \lambda_1(A) + \delta\mu$, wobei $|\delta\mu| \leq \|B\|^2/(g_{N_0+1} - \|B\|)$ und $B$ die off-diagonale Kopplung ist. Die Neumann-Reihe $(I + R_0 K)^{-1} = \sum_{n \geq 0} (-R_0 K)^n$ konvergiert, wenn $\|R_0 K\| < 1$, und liefert die Trunkierungsfehler-Schranke. Die numerischen Werte werden aus den Galerkin-Matrizen an jedem Zertifikatwert berechnet. Der $\cos$-Sektor hat $\|R_0 K\| \approx 0.50$ (stabil, da der 4-Moden-Block die dominante Mode erfasst), während der $\sin$-Sektor größere $\lambda$ ($\geq 500$) für die Konvergenz benötigt. Die Korrektur der \emph{Differenz} ist beschränkt durch $\|R_0 K\|^2 \cdot |E_{\sin} - E_{\cos}| = O(\|R_0 K\|^2)$, da $|E_{\sin} - E_{\cos}| = O(1)$ (\Cref{lem:leading-cancel}). Division durch $D = \Theta(\sqrt{\lambda})$ ergibt $o(1)$.
\end{proof}


% =============================================================================
\section{Verbindungen zu existierenden Arbeiten}
% =============================================================================

\subsection{Li-Keiper-Koeffizienten}
Die Koeffizienten $\lambda_n$ (auch Li-Keiper-Koeffizienten genannt) wurden von Keiper \cite{Keiper1992} numerisch berechnet, von Li \cite{Li1997} als Kriterium bewiesen und von Bombieri--Lagarias \cite{BombieriLagarias1999} in der hier verwendeten Darstellung etabliert.

\subsection{Weils explizite Formel und Connes' Programm}
Die explizite Formel von Weil
\[ \sum_\rho h(\rho) = h(0) + h(1) - \sum_p \sum_{m=1}^\infty \frac{\log p}{p^{m/2}}\,\hat{h}(m\log p) \]
stellt eine direkte Brücke zwischen Nullstellen und Primzahlen her. Connes \cite{Connes2025} konstruiert daraus einen selbstadjungierten Operator $A_\lambda$ und reduziert RH auf die Einfachheit seines Grundzustands mit gerader Eigenfunktion (Theorem~6.1). Die gerade Dominanz ist eine \emph{unbewiesene Annahme} in allen aktuellen Ansätzen. Unsere Zerlegungsanalyse zeigt, dass die Primzahlbeiträge, nicht der archimedische Kern, der einzige Treiber der geraden Dominanz sind.


% =============================================================================
\section{Schlussfolgerung}
% =============================================================================

Diese Arbeit etabliert die gerade Dominanz der quadratischen Weil-Form $\mathrm{QW}_\lambda$ für alle $\lambda \geq 100$. Der Beweis kombiniert drei Zutaten: (1)~das Shift-Parity-Lemma, das zeigt, dass jede Primzahl individuell gerade Eigenfunktionen bevorzugt; (2)~$33$ computergestützte Zertifikate, die eine rigorose Verifikation mittels Intervallarithmetik liefern; und (3)~den M1$''$-Resolventen-Rahmen mit PNT-Transfer, der die gerade Dominanz analytisch für alle hinreichend großen $\lambda$ etabliert. Das Drei-Regime-Brückenargument (Proposition~A6) vereint diese Zutaten zu einem lückenlosen Beweis für alle $\lambda \geq 100$.

Die zentrale analytische Einsicht ist das Leading-Mode-Cancellation-Lemma: Die Kopplungsdifferenz zwischen geradem und ungeradem Sektor ist durch $O(1)$ beschränkt, während der Diagonalvorteil $D(\lambda)$ wie $\sqrt{\lambda}$ wächst, was ein Verhältnis $\rho(\lambda) \to 0$ erzeugt. Kombiniert mit der Euler--Maclaurin-Proposition ($\rho^{\mathrm{EM}} < 0$ für $L \geq 14$) und dem PNT-Transfer-Lemma ergibt sich M1$''$ (Resolventen-Subdominanz) für alle $\lambda \geq \lambda_0$.

Zwei numerische Inputs verbleiben: die spektrale Gap-Konstante $c_{\mathrm{gap}} \geq 0.5$ (\Cref{lem:higher-mode-decay}) und die Neumann-Konvergenzbedingung $\|R_0 K\| < 1$ (\Cref{lem:resolvent-truncation}), beide an allen Zertifikatwerten verifiziert. Analytische Herleitungen würden einen vollständig analytischen Beweis für Regime~2 vervollständigen.


\section*{KI-Offenlegung}

Dieses Manuskript wurde in umfassender Zusammenarbeit mit KI-Sprachmodellen (Claude, Anthropic) erstellt. Die KI trug zur konzeptionellen Exploration, analytischen Arbeit, Literaturrecherche, Texterstellung und kritischen Überprüfung bei. Alle computergestützten Zertifikate wurden durch deterministische numerische Skripte unabhängig von der KI berechnet. Der Autor, Lukas Geiger, hat die Forschungsrichtung konzipiert, seine Fachexpertise eingebracht und die finale redaktionelle Entscheidungsgewalt ausgeübt. Der Autor übernimmt die alleinige wissenschaftliche Verantwortung für den gesamten Inhalt einschließlich etwaiger Fehler.

\begin{thebibliography}{99}

\bibitem{Geiger2026partI}
L.~Geiger, \emph{Das FST-RH Framework: Grundlagen, Hindernisse und Neuausrichtung}, 2026.

\bibitem{Geiger2026partIII}
L.~Geiger, \emph{Das FST-RH Framework: Conclusio}, 2026.

\bibitem{Connes2025}
A.~Connes, \emph{Prolate and the Riemann zeta function}, arXiv:2602.04022v1, 2026.

\bibitem{ConnesvS2025}
A.~Connes und W.~D. van Suijlekom, \emph{Spectral truncation of the Riemann zeta function and a Carath\'eodory--Fej\'er approximation}, Preprint, 2025.

\bibitem{BombieriLagarias1999}
E.~Bombieri und J.~C.~Lagarias, \emph{Complements to Li's criterion for the Riemann hypothesis}, J.~Number Theory~\textbf{77} (1999), 274--287.

\bibitem{Li1997}
X.-J.~Li, \emph{The positivity of a sequence of numbers and the Riemann hypothesis}, J.~Number Theory~\textbf{65} (1997), 325--333.

\bibitem{Keiper1992}
J.~B.~Keiper, \emph{Power series expansions of Riemann's $\xi$ function}, Math.~Comp.~\textbf{58} (1992), 765--773.

\bibitem{Connes1999}
A.~Connes, \emph{Trace formula in noncommutative geometry and the zeros of the Riemann zeta function}, Selecta Math.~\textbf{5} (1999), 29--106.

\bibitem{Deninger2002}
C.~Deninger, \emph{On the nature of the ``explicit formulas'' in analytic number theory}, Banach Center Publ.~\textbf{57} (2002), 97--118.

\bibitem{IwaniecKowalski2004}
H.~Iwaniec und E.~Kowalski, \emph{Analytic Number Theory}, AMS Colloquium Publ.~\textbf{53}, 2004.

\bibitem{Kato1995}
T.~Kato, \emph{Perturbation Theory for Linear Operators}, Springer, 1995.

\end{thebibliography}

\end{document}
